\chapter{Amortized Analysis}

Across our study of dynamic programming and greedy choices, we analyzed algorithms where each step was bounded by a uniform worst-case cost. But when designing dynamic data structures, a different pattern emerges. You build a data
structure. You measure one operation on it and find that, every so often, that
operation is \emph{slow} -- it does a big burst of work, touching a large chunk
of the structure. Naively, you write down ``worst case per operation: expensive''
and conclude the whole sequence of operations must be expensive. And yet, when
you actually run a long sequence of these operations and divide the total time by
the number of operations, the average comes out \emph{cheap} -- often constant.
The slow operations are real, but they are so rare, and the cheap ones so common,
that the expensive bursts are paid for by everything around them.

\emph{Amortized analysis} is the mathematical framework for making that intuition precise. It
lets you charge a sequence of operations fairly -- not by pretending the slow
ones do not happen, but by spreading their cost over the many cheap operations
that surround them. The payoff is that you can prove a tight bound on
the \emph{total} cost of a sequence without ever invoking probability. There is no
``on average over random inputs'' here. The bound holds for \emph{every} sequence,
in the worst case, which is what makes it a guarantee you can actually rely on.

We will meet three techniques -- the
\emph{aggregate}, \emph{accounting}, and \emph{potential} methods -- and run each
of them on the same three examples: a stack with a bulk-pop operation, a binary counter you increment, and a
dynamic table that doubles itself when it fills up. By the end you will be able to
look at a bursty data structure and say, calmly, ``yes, but amortized it's
constant time, and here is the proof.''

\chapterline

\section{Worst case per operation versus worst case per sequence}

Let us be very careful about what question we are asking, because the entire
subject rests on this distinction.

When we analyzed sorting, we asked: what is the worst-case running time of
\emph{one call} to this algorithm on an input of size $n$? That is a perfectly
good question. But data structures are different. You rarely call a data structure
once. You hammer it with a long \emph{sequence} of operations -- push, push, pop,
push, search, insert, delete -- and what you actually care about is the total time
for the whole sequence.

\begin{intuitionbox}
Think about your monthly expenses. Most days you spend almost nothing. But once a
month the rent comes due, and that single day is enormous. If a visitor watched
you on rent day, they would scream ``you can't afford this lifestyle!'' But you
can, because you have all those cheap days to balance it out. Your \emph{average
daily spend} is modest. Amortized cost is exactly this average daily spend,
computed across the whole month -- and computed so that it is a true
\emph{upper bound} no matter which days the visitor happens to look at.
\end{intuitionbox}

Here is the key definition, stated plainly.

\begin{definition}[Amortized cost]
\label{ch17:def-amortized}
Suppose a sequence of $n$ operations runs on a data structure and the total actual
cost of the whole sequence is $T(n)$. We say each operation has \emph{amortized
cost} at most $\hat{c}$ if, for \emph{every} such sequence,
\[
  T(n) \;\le\; \hat{c} \cdot n.
\]
More generally we may assign different amortized costs $\hat{c}_i$ to different
operations $i$, and the requirement is that the sum of amortized costs is an upper
bound on the total actual cost:
$\sum_{i=1}^{n} c_i \le \sum_{i=1}^{n} \hat{c}_i$, where $c_i$ is the actual cost
of the $i$th operation.
\end{definition}

That last inequality is the core of the method. The
amortized costs are \emph{bookkeeping numbers we invent}. We are free to choose
them however we like, as long as one rule holds: their running sum must always be
at least the running sum of the real costs. If we can do that, then the total
amortized cost is a genuine upper bound on the total real cost, and we have
permission to use the amortized cost as if it were the per-operation cost.

\begin{focusbox}
Amortized cost is not a measured quantity. It is an \emph{accounting choice} that
you must \emph{justify}. The justification always reduces to one inequality: across
any prefix of the sequence, total amortized cost $\ge$ total actual cost. Get that
inequality and you may quote the amortized cost freely.
\end{focusbox}

\begin{trapbox}
Amortized is \emph{not} average-case, and it is \emph{not} expected. There is no
randomness anywhere in this chapter. An amortized bound is a worst-case guarantee
about a \emph{sequence}; an average-case bound is a probabilistic statement about a
\emph{single} operation over random inputs. A randomized algorithm can have a great
expected cost and still, on one unlucky run, blow up. An amortized bound never
blows up over the sequence: it holds for every sequence, every time, with no luck
involved.
\end{trapbox}

We will now develop three methods that all establish bounds of this form. They
prove the same kinds of theorems but use three different \emph{styles of
argument}, and different examples feel natural in
different styles. A good engineer keeps all three in the toolbox.

\section{The aggregate method: just add it all up}

The aggregate method is the bluntest of the three. You do not try
to be clever about individual operations. You ask one question: across the whole
sequence of $n$ operations, what is the total work $T(n)$? Then you divide by $n$
and call the result the amortized cost per operation. That is it.

\begin{intuitionbox}
The aggregate method is the ``total the receipt, then split the bill evenly''
approach. You do not care who ordered the lobster and who ordered the salad. You
sum the entire tab and divide by the number of diners. If the total is small, the
per-person share is small, even if one dish was expensive.
\end{intuitionbox}

It sounds too simple to be useful, but for many structures the total $T(n)$ has a
clean closed form that is much smaller than (worst-case-per-operation) $\times n$.
Let us see this on our first running example.

\subsection{Example: the stack with \textsc{Multipop}}

Start with an ordinary stack. It supports two cheap operations:
\textsc{Push}($S, x$) puts $x$ on top, and \textsc{Pop}($S$) removes and returns
the top. Each costs $1$ unit (one object moved). Now we add a third, bulkier
operation.

\begin{algorithm}[H]
\caption{\textsc{Multipop}($S, k$)}
\KwIn{A stack $S$ and a count $k \ge 0$.}
\KwOut{The top $\min(k, |S|)$ elements are popped off.}
\While{$S$ is not empty \textbf{and} $k > 0$}{
  $\textsc{Pop}(S)$\;
  $k \leftarrow k - 1$\;
}
\end{algorithm}

\textsc{Multipop} just pops repeatedly until either it has popped $k$ items or the
stack runs dry. Its actual cost is the number of \textsc{Pop} calls it makes,
which is $\min(k, |S|)$.

Now here is the trap a naive analysis falls into. A single \textsc{Multipop} can
be expensive: if the stack holds $n$ items and you call \textsc{Multipop}($S, n$),
it does $n$ units of work. So the worst-case cost of one operation is $\Theta(n)$.
A sequence of $n$ operations, naively bounded, looks like it could cost
$n \cdot \Theta(n) = \Theta(n^2)$.

\begin{trapbox}
That $\Theta(n^2)$ bound is correct -- it is a valid upper bound -- but it is
wildly loose. It assumes \emph{every} operation can be the expensive one. It cannot.
A \textsc{Multipop} that pops $100$ items requires that $100$ items were pushed
earlier, and pushing them already cost $100$ units. You cannot pop what you never
pushed. The naive bound double-counts a constraint that the structure enforces for
free.
\end{trapbox}

Let us total it up properly. Consider any sequence of $n$ operations, each one a
\textsc{Push}, \textsc{Pop}, or \textsc{Multipop}, starting from an empty stack.
The decisive observation is this:

\begin{focusbox}
Every object can be popped \emph{at most once for each time it is pushed}. The
total number of pops (counting each pop inside a \textsc{Multipop} individually)
can therefore never exceed the total number of pushes.
\end{focusbox}

How many pushes can there be in a sequence of $n$ operations? At most $n$, since
each operation is a single \textsc{Push} at most. So the total number of pushes is
$\le n$, and therefore the total number of pops -- across all \textsc{Pop} and
\textsc{Multipop} calls combined -- is also $\le n$. Each push costs $1$ and each
pop costs $1$, so:
\[
  T(n) \;=\; (\text{total pushes}) + (\text{total pops}) \;\le\; n + n \;=\; 2n
  \;=\; O(n).
\]
Divide by $n$: the amortized cost per operation is at most $2n / n = 2 = O(1)$.

\begin{aestheticbox}
Notice how much the bound improves: from a naive $\Theta(n^2)$ down to a tight
$\Theta(n)$ total, just by counting pushes and pops globally instead of bounding
each operation separately. The expensive \textsc{Multipop}s are real, but they are
\emph{funded} by the pushes that came before them. This same ``you can only take
out what you put in'' accounting underlies the analysis of countless real systems:
memory allocators, log compaction, garbage collectors. Whenever work is bounded by
some resource that had to be created earlier, amortized analysis usually applies.
\end{aestheticbox}

\begin{labbox}
Trace it yourself. Start with an empty stack and run this exact sequence of $8$
operations, writing down the stack and the actual cost after each:
\textsc{Push}(a), \textsc{Push}(b), \textsc{Push}(c),
\textsc{Multipop}($2$), \textsc{Push}(d), \textsc{Push}(e),
\textsc{Multipop}($5$), \textsc{Push}(f).
Add the actual costs (you should get $3$ pushes then a $2$-pop, $2$ pushes, a
$3$-pop because only $3$ items are present, then $1$ push: $1+1+1+2+1+1+3+1 = 11$).
Confirm $11 \le 2 \cdot 8 = 16$. Notice the $5$-multipop only cost $3$: it could
not pop what was not there.
\end{labbox}

\subsection{Example: the binary counter}

Our second running example will follow us through all three methods, so meet it
now. We keep a $k$-bit binary counter in an array $A[0 \mathinner{\ldotp\ldotp} k-1]$, where $A[0]$
is the lowest-order bit. The only operation is \textsc{Increment}, which adds $1$.

\begin{algorithm}[H]
\caption{\textsc{Increment}($A$)}
\KwIn{A binary counter stored in bits $A[0 \mathinner{\ldotp\ldotp} k-1]$, low bit first.}
\KwOut{$A$ holds the value increased by $1$ (modulo $2^k$).}
$i \leftarrow 0$\;
\While{$i < k$ \textbf{and} $A[i] = 1$}{
  $A[i] \leftarrow 0$ \tcp*{a carry: flip a 1 down to 0}
  $i \leftarrow i + 1$\;
}
\If{$i < k$}{
  $A[i] \leftarrow 1$ \tcp*{set the lowest 0 bit to 1}
}
\end{algorithm}

This is just grade-school carrying in base two. To add $1$, you flip the trailing
run of $1$s to $0$s, then set the next $0$ to a $1$. The cost of one
\textsc{Increment} is the number of bits it flips.

How bad can one \textsc{Increment} be? When the counter reads $0111\cdots1$ (all
low bits set), incrementing flips \emph{all} of them: up to $k$ bit-flips. So the
worst case of a single operation is $\Theta(k)$, and a naive bound on $n$
increments is $n \cdot \Theta(k) = \Theta(nk)$. Let us see how loose that is.

The aggregate insight is to ask not ``how many bits does one increment flip?'' but
``how many times, total, does each \emph{individual bit} flip across $n$
increments?'' Watch the bits.

\begin{itemize}
  \item Bit $A[0]$ flips on \emph{every} increment: $1,0,1,0,\ldots$. So it flips
        $n$ times.
  \item Bit $A[1]$ flips on every \emph{second} increment: it changes once for
        every two ticks of $A[0]$. So it flips $\lfloor n/2 \rfloor$ times.
  \item Bit $A[2]$ flips every \emph{fourth} increment: $\lfloor n/4 \rfloor$
        times.
  \item In general, bit $A[i]$ flips $\lfloor n / 2^i \rfloor$ times.
\end{itemize}

Figure~\ref{ch17:counter-flips} makes this halving-per-column pattern visible for a
small counter.

\begin{intuitionbox}
A binary counter is a row of toggle switches geared like an odometer. The ones
digit spins fast and flips constantly. The tens digit barely moves. The
highest-order switch flips only on the rare grand rollover. Most of the bits
rarely move, so the total flipping work is dominated by the low bits, and each of
their flips is \emph{cheap}.
\end{intuitionbox}

\begin{figure}[h]
\centering
\begin{tikzpicture}[
    cell/.style={draw,minimum width=0.55cm,minimum height=0.45cm,font=\ttfamily\footnotesize},
    flip/.style={cell,fill=quotientorange!45},
    same/.style={cell,fill=gray!8}]
  % Columns left->right are bits A3 A2 A1 A0. Each entry: row/col/bit/flag
  % flag f = flipped on this increment (shaded), s = unchanged.
  \foreach \r/\k/\b/\f in {
    0/0/0/s, 0/1/0/s, 0/2/0/s, 0/3/0/s,
    1/0/0/s, 1/1/0/s, 1/2/0/s, 1/3/1/f,
    2/0/0/s, 2/1/0/s, 2/2/1/f, 2/3/0/f,
    3/0/0/s, 3/1/0/s, 3/2/1/s, 3/3/1/f,
    4/0/0/s, 4/1/1/f, 4/2/0/f, 4/3/0/f,
    5/0/0/s, 5/1/1/s, 5/2/0/s, 5/3/1/f,
    6/0/0/s, 6/1/1/s, 6/2/1/f, 6/3/0/f,
    7/0/0/s, 7/1/1/s, 7/2/1/s, 7/3/1/f,
    8/0/1/f, 8/1/0/f, 8/2/0/f, 8/3/0/f}{
    \ifnum\pdfstrcmp{\f}{f}=0\relax
      \node[flip] at (0.55*\k,-0.5*\r) {\b};
    \else
      \node[same] at (0.55*\k,-0.5*\r) {\b};
    \fi
  }
  \foreach \r in {0,...,8}{
    \node[font=\rmfamily\footnotesize] at (-1.0,-0.5*\r) {\r};
  }
  \node[font=\rmfamily\footnotesize] at (-1.0,0.55) {value};
  \foreach \k/\lbl in {0/$A_3$,1/$A_2$,2/$A_1$,3/$A_0$}{
    \node[font=\rmfamily\footnotesize] at (0.55*\k,0.55) {\lbl};
  }
\end{tikzpicture}
\caption{A $4$-bit counter climbing $0 \to 8$. Shaded cells are the bits that
flipped on that increment. The low bit $A_0$ flips every step; $A_1$ every two
steps; $A_2$ every four; $A_3$ only on the big rollover at $7 \to 8$. Each column
flips half as often as the one to its right, which gives the geometric series
$n + n/2 + n/4 + \cdots < 2n$.}
\label{ch17:counter-flips}
\end{figure}

Now add up the total number of flips across all bits:
\[
  T(n) \;=\; \sum_{i=0}^{k-1} \left\lfloor \frac{n}{2^i} \right\rfloor
  \;<\; \sum_{i=0}^{\infty} \frac{n}{2^i}
  \;=\; n \sum_{i=0}^{\infty} \frac{1}{2^i}
  \;=\; n \cdot 2
  \;=\; 2n.
\]
That geometric series $1 + \tfrac12 + \tfrac14 + \cdots = 2$ gives the bound.
The total flip count over $n$ increments is less than $2n$, so:
\[
  T(n) = O(n), \qquad \text{amortized cost per \textsc{Increment}} = \frac{T(n)}{n} < 2 = O(1).
\]

\begin{aestheticbox}
The naive bound was $\Theta(nk)$; the truth is $\Theta(n)$. The factor of $k$
disappears, because the high bits almost never move. This is the same
reason that incrementing a $64$-bit loop counter is effectively free even though,
once in a great while, an increment ripples through dozens of carries. Because
flip frequencies halve with each higher bit, the geometric series ensures
that all the rare carries combined cost no more than the common low-bit flips.
\end{aestheticbox}

\begin{labbox}
Build the table by hand. Start at $0000$ and increment a $4$-bit counter $16$
times back to $0000$. After each increment, write the counter value and the number
of bits that flipped. You should see the cost pattern
$1,2,1,3,1,2,1,4,1,2,1,3,1,2,1,4$. Sum it: you get $30$. Check that this is the
total flips $\lfloor 16/1\rfloor + \lfloor 16/2\rfloor + \lfloor 16/4\rfloor +
\lfloor 16/8\rfloor = 16+8+4+2 = 30$, and that $30 < 2 \cdot 16 = 32$. Notice the
``$4$'' costs appear only at the rollovers $7\to8$ and $15\to0$ -- exactly twice.
\end{labbox}

The aggregate method works well when a clean global total exists. But sometimes
you want to assign \emph{different} amortized costs to different operations, or you
want a more local, more flexible argument. That is where the next two methods come
in.

\section{The accounting method: prepay and bank the credit}

The accounting method uses a bookkeeping device: we charge each operation an
amortized cost that may differ from its real cost, and we treat the difference as
\emph{currency}. If we charge an operation more than it actually costs, the surplus
is stored in the data structure as \emph{credit}. If an operation costs more than
we charge it, it pays the difference by spending stored credit.

\begin{intuitionbox}
Imagine each object in your structure carries a little prepaid debit card. When
you do a cheap operation, you overcharge it slightly and load the extra onto the
card. Later, when an expensive operation happens, it does not reach into your
wallet -- it swipes the cards that were prepaid earlier. As long as no card ever
goes negative, you know the prepaid charges covered all the real work. The total
you charged is then a valid upper bound on the total real cost.
\end{intuitionbox}

The single rule that makes this rigorous:

\begin{focusbox}
The bank balance -- total credit stored in the structure -- must \emph{never go
negative}. Equivalently, at every prefix of the sequence, the total amount charged
(amortized) must be at least the total real cost. If credit can never go negative,
then total charged $\ge$ total spent, which is exactly
Definition~\ref{ch17:def-amortized}. The credit you store represents prepaid future
work, so it makes physical sense that it cannot be negative -- you cannot prepay a
debt you have not funded.
\end{focusbox}

The hard part of the accounting method is \emph{choosing the charges}. You want simple,
constant charges that nevertheless keep the bank solvent. Let us pick charges for
our two examples.

\subsection{Accounting for the stack}

We assign amortized charges as follows:
\begin{itemize}
  \item \textsc{Push}: charge $2$. Real cost is $1$. The extra $1$ stays as credit
        \emph{on the pushed object}.
  \item \textsc{Pop}: charge $0$. Real cost is $1$. Pay the $1$ with the credit
        sitting on the object being popped.
  \item \textsc{Multipop}: charge $0$. Real cost is the number of items it pops;
        each pop is paid by the credit on the object it removes.
\end{itemize}

Here is the invariant that makes the scheme work: \emph{every object currently on the
stack carries exactly $\$1$ of credit}. Why is this maintained? When you push an
object, you charged $\$2$: $\$1$ pays the real push cost, and $\$1$ rides onto the
object. So freshly pushed objects satisfy the invariant. When you pop an object
(whether via \textsc{Pop} or inside \textsc{Multipop}), the object carries exactly
the $\$1$ needed to pay for its own removal. The credit is consumed precisely when
the object leaves. Figure~\ref{ch17:stack-credit} shows the coins in action across
a push burst followed by a \textsc{Multipop}.

\begin{intuitionbox}
Every object walks onto the stack already holding the coin it will need to walk
back off. \textsc{Push} hands it that coin. Whenever it is popped -- now, later, in
bulk, it does not matter -- it pays its own way out. No pop ever finds an empty
pocket, so the bank never goes negative.
\end{intuitionbox}

\begin{figure}[h]
\centering
\begin{tikzpicture}[font=\small,
    item/.style={draw,minimum width=1.1cm,minimum height=0.55cm,fill=sourceblue!12},
    coin/.style={draw,circle,inner sep=0pt,minimum size=0.42cm,fill=quotientorange!40,font=\tiny}]
  % State after Push a, Push b, Push c (each carrying a $1 coin)
  \node at (1.5,3.1) {after \textsc{Push}(a,b,c)};
  \node[item] (a) at (1.5,0.0) {a}; \node[coin,right=2pt of a] {\$1};
  \node[item] (b) at (1.5,0.6) {b}; \node[coin,right=2pt of b] {\$1};
  \node[item] (c) at (1.5,1.2) {c}; \node[coin,right=2pt of c] {\$1};
  \draw[thick] (0.8,-0.3) -- (0.8,2.0) (0.8,-0.3) -- (2.2,-0.3) (2.2,-0.3) -- (2.2,2.0);
  % arrow
  \draw[-{Stealth[length=3mm]},structurepurple,thick] (3.2,0.8) -- (4.6,0.8)
    node[midway,above,font=\tiny] {\textsc{Multipop}(2)};
  % State after Multipop(2): c and b popped, paying their own coins
  \node at (6.3,3.1) {after \textsc{Multipop}(2)};
  \node[item] (a2) at (6.3,0.0) {a}; \node[coin,right=2pt of a2] {\$1};
  \draw[thick] (5.6,-0.3) -- (5.6,2.0) (5.6,-0.3) -- (7.0,-0.3) (7.0,-0.3) -- (7.0,2.0);
  \node[font=\tiny,text=trapred!90!black,align=center] at (6.3,2.45)
    {$b,c$ popped\\(each spent its \$1)};
\end{tikzpicture}
\caption{The accounting view of the stack. Each \textsc{Push} overcharges by \$1
and parks that coin on the object (left). When \textsc{Multipop}(2) removes
$c$ and $b$, each pays for its own removal out of the coin it has been carrying
(right). The bank -- the total of all parked coins -- equals the number of items
on the stack and never dips below zero.}
\label{ch17:stack-credit}
\end{figure}

Since the bank never goes negative, the total charged is an upper bound on total
real cost. The amortized cost of each operation is at most $2$ (the most we ever
charge), so $n$ operations cost at most $2n = O(n)$. Same bound as the aggregate
method, reached by a completely different argument -- and notice this argument is
more \emph{local}: it never needed to know the length of the sequence.

\subsection{Accounting for the binary counter}

For the counter, the trick is to ask: which bit-flips are expensive, and which can
prepay them? An \textsc{Increment} flips a run of $1$s down to $0$, then sets one
$0$ up to $1$. The flips $1 \to 0$ can be arbitrarily numerous (that is the carry
chain), but a $0 \to 1$ flip -- setting a bit -- happens \emph{exactly once} per
\textsc{Increment}. That asymmetry is what makes the accounting work.

We charge:
\begin{itemize}
  \item Setting a bit ($0 \to 1$): charge $\$2$. One dollar pays the real flip;
        the other dollar rides onto the bit as credit.
  \item Clearing a bit ($1 \to 0$): charge $\$0$. Pay its real flip with the $\$1$
        credit that the bit has been carrying since it was set.
\end{itemize}

Invariant: \emph{every bit that currently equals $1$ carries exactly $\$1$ of
credit}. When we set a bit to $1$, we loaded $\$1$ onto it. When a carry later
clears that bit to $0$, the stored dollar pays for the clear. Because every
\textsc{Increment} sets at most one bit to $1$, every \textsc{Increment} is charged
at most $\$2$. Therefore $n$ increments cost at most $2n = O(n)$, and the amortized
cost per \textsc{Increment} is at most $2 = O(1)$.

\begin{figure}[h]
\centering
\begin{tikzpicture}[font=\small]
  % Lane 1: Before increment (Value 11 = 1011_2)
  \begin{scope}
    \node[font=\footnotesize\scshape, text=gray!70!black, anchor=west] at (-0.5, 1.6) {Before: Counter $= 11$ ($(1011)_2$), $3$ coins stored on the $1$-bits};
    \foreach \b/\val/\hascoin in {3/1/1, 2/0/0, 1/1/1, 0/1/1} {
      \draw[thick, draw=gray!60, fill=gray!8, rounded corners=2pt] ({(3-\b)*1.4}, 0) rectangle ({(3-\b)*1.4+1.1}, 0.75);
      \node[font=\small\ttfamily\bfseries] at ({(3-\b)*1.4+0.55}, 0.375) {\val};
      \node[font=\tiny, text=gray!60] at ({(3-\b)*1.4+0.55}, -0.3) {$A[\b]$};
      \ifnum\hascoin=1
        \node[draw=quotientorange, fill=quotientorange!35, circle, thick, inner sep=1.5pt, minimum size=13pt, font=\tiny\bfseries] 
          at ({(3-\b)*1.4+0.55}, 1.05) {\$1};
      \fi
    }
  \end{scope}

  % Transition Arrow
  \draw[-{Latex[length=2.5mm]}, very thick, quotientorange] (2.2, -0.6) -- 
    node[right=3pt, font=\scriptsize, text=quotientorange] {\textsc{Increment}: amortized charge $\$2$} (2.2, -1.6);

  % Lane 2: After increment (Value 12 = 1100_2)
  \begin{scope}[yshift=-3.8cm]
    \node[font=\footnotesize\scshape, text=gray!70!black, anchor=west] at (-0.5, 1.6) {After: Counter $= 12$ ($(1100)_2$), bits $0,1$ paid with credit; bit $2$ gained $\$1$};
    \foreach \b/\val/\hascoin/\status in {3/1/1/{same}, 2/1/1/{new}, 1/0/0/{spent}, 0/0/0/{spent}} {
      \ifnum\b=2
        \draw[thick, draw=targetgreen, fill=targetgreen!20, rounded corners=2pt] ({(3-\b)*1.4}, 0) rectangle ({(3-\b)*1.4+1.1}, 0.75);
      \else
        \draw[thick, draw=gray!60, fill=gray!8, rounded corners=2pt] ({(3-\b)*1.4}, 0) rectangle ({(3-\b)*1.4+1.1}, 0.75);
      \fi
      \node[font=\small\ttfamily\bfseries] at ({(3-\b)*1.4+0.55}, 0.375) {\val};
      \node[font=\tiny, text=gray!60] at ({(3-\b)*1.4+0.55}, -0.3) {$A[\b]$};
      \ifnum\hascoin=1
        \node[draw=quotientorange, fill=quotientorange!35, circle, thick, inner sep=1.5pt, minimum size=13pt, font=\tiny\bfseries] 
          at ({(3-\b)*1.4+0.55}, 1.05) {\$1};
      \fi
    }
    % Annotation on spent coins
    \node[font=\tiny, text=trapred!80!black] at ({(3-0)*1.4+0.55}, 1.05) {spent};
    \node[font=\tiny, text=trapred!80!black] at ({(3-1)*1.4+0.55}, 1.05) {spent};
    \node[font=\tiny\bfseries, text=targetgreen!90!black, right=1pt] at ({(3-2)*1.4+0.85}, 1.05) {+\$1};
  \end{scope}
\end{tikzpicture}
\caption{The accounting method on a binary counter. Every bit equal to $1$ maintains exactly $\$1$ of stored credit. Setting a bit ($0 \to 1$) incurs an amortized charge of $\$2$: $\$1$ pays for the physical flip, and $\$1$ is banked on the bit. Clearing a bit ($1 \to 0$) is charged $\$0$, paid for by releasing the banked dollar. The bank balance equals the number of $1$-bits, preventing the bank from ever going negative.}
\label{ch17:fig-counter-accounting}
\end{figure}

\begin{focusbox}
The same pattern appears in both: charge $\$2$ for the one cheap, guaranteed,
once-per-operation action (a push, a bit-set), and let the prepaid dollar fund the
unbounded-looking burst later (a multipop, a carry chain). The structure of the
argument is identical even though the structures are different.
\end{focusbox}

\begin{walkbox}
Take this on a walk. In the counter argument we charged the \emph{set} and let the
\emph{clears} be free. Could we instead charge each clear and make sets free?
Try to define such a scheme and find where the bank goes negative. What is special
about ``at most one set per increment'' that makes the set the right thing to
charge? When you see why one direction works and the other fails, you understand
the accounting method.
\end{walkbox}

\section{The potential method: bottle the credit into a single number}

The accounting method spread credit across individual objects -- a dollar on each
stacked item, a dollar on each set bit. The \emph{potential method} is the same
idea, viewed globally. Instead of tracking who holds which coin, we summarize
the \emph{entire} stored credit as one number: the \emph{potential} $\Phi$ of the
whole data structure. As the structure evolves, $\Phi$ rises and falls, and that
rise and fall is exactly the bookkeeping.

\begin{intuitionbox}
Think of potential energy in physics. Lift a weight and you store energy; the
storing costs you effort now. Let it drop and that stored energy does work for
free later. The data structure's potential $\Phi$ is precisely this: a measure of
``stored-up future work.'' Cheap operations that build up structure raise $\Phi$
(you pay a little extra now); the rare expensive operation that tears that
structure down releases $\Phi$ (it pays for itself out of the reservoir).
\end{intuitionbox}

Here is the precise machinery. Let $D_0$ be the starting state and $D_i$ the state
after the $i$th operation. We pick a \emph{potential function} $\Phi$ mapping each
state to a real number. The amortized cost of the $i$th operation is \emph{defined}
as its actual cost plus the change in potential it causes:
\[
  \boxed{\;\hat{c}_i \;=\; c_i \;+\; \Phi(D_i) - \Phi(D_{i-1}).\;}
\]

Why is this the right definition? Sum the amortized costs over the whole sequence
and watch the potentials telescope:
\[
  \sum_{i=1}^{n} \hat{c}_i
  = \sum_{i=1}^{n}\bigl(c_i + \Phi(D_i) - \Phi(D_{i-1})\bigr)
  = \left(\sum_{i=1}^{n} c_i\right) + \Phi(D_n) - \Phi(D_0).
\]
Every intermediate $\Phi(D_i)$ cancels: $+\Phi(D_1)$ from the first term against
$-\Phi(D_1)$ from the second, and so on. Only the endpoints survive. Rearranging,
the total \emph{actual} cost is
\[
  \sum_{i=1}^{n} c_i \;=\; \sum_{i=1}^{n}\hat{c}_i \;-\; \Phi(D_n) + \Phi(D_0).
\]

\begin{formalbox}
If we choose $\Phi$ so that $\Phi(D_0) = 0$ and $\Phi(D_i) \ge 0$ for all $i$,
then $-\Phi(D_n) + \Phi(D_0) = -\Phi(D_n) \le 0$, and therefore
\[
  \sum_{i=1}^{n} c_i \;\le\; \sum_{i=1}^{n} \hat{c}_i.
\]
So the total amortized cost is an upper bound on the total actual cost -- exactly
the requirement of Definition~\ref{ch17:def-amortized}. The two conditions
``$\Phi$ starts at $0$'' and ``$\Phi$ never goes negative'' are the potential-method
analogue of ``the bank never goes negative.''
\end{formalbox}

\begin{focusbox}
The potential is just the bank balance from the accounting method, packaged as a
single function of the state. Conditions $\Phi(D_0)=0$ and $\Phi \ge 0$ say exactly
``start with no credit, and never go into debt.'' In other words, the potential
method is the accounting method with the entire credit ledger summed into a
single number.
\end{focusbox}

\subsection{Potential for the stack}

Define $\Phi(D) = $ the number of objects currently on the stack. Initially the
stack is empty so $\Phi(D_0) = 0$, and a count of objects is never negative, so
$\Phi \ge 0$ always. Both conditions hold. Now compute the amortized cost of each
operation as ``actual cost plus change in $\Phi$.''

\textbf{Push.} Actual cost $1$. The stack grows by one object, so $\Phi$ goes up by
$1$. Amortized cost $= 1 + 1 = 2$.

\textbf{Pop.} Actual cost $1$. The stack shrinks by one, so $\Phi$ drops by $1$.
Amortized cost $= 1 + (-1) = 0$.

\textbf{Multipop} popping $m = \min(k, |S|)$ items. Actual cost $m$. The stack
shrinks by $m$, so $\Phi$ drops by $m$. Amortized cost $= m + (-m) = 0$.

Every operation has amortized cost at most $2$. So $n$ operations cost at most
$2n = O(n)$. Notice this reproduces the accounting result -- ``$\$1$ on each
stacked object'' is literally ``$\Phi = $ number of objects.'' The potential method
just did the same accounting without ever naming an individual coin.

\subsection{Potential for the binary counter}

Define $\Phi(D) = b_i$, the number of $1$-bits currently in the counter after
operation $i$. Start at $0$, so $\Phi(D_0) = 0$, and a bit count is never negative,
so $\Phi \ge 0$. Both conditions hold.

Now analyze a single \textsc{Increment}. Suppose it clears a run of $t$ trailing
$1$-bits to $0$ (those are the $t$ carries) and then sets one $0$-bit to $1$ --
unless the counter overflowed, but let us take the typical, non-overflow case
first.

\begin{itemize}
  \item \textbf{Actual cost} $c_i = t + 1$: it flipped $t$ ones down and one zero
        up.
  \item \textbf{Change in potential}: we lost $t$ ones (the cleared bits) and
        gained $1$ one (the set bit), so $b_i = b_{i-1} - t + 1$, giving
        $\Phi(D_i) - \Phi(D_{i-1}) = (b_{i-1} - t + 1) - b_{i-1} = 1 - t$.
  \item \textbf{Amortized cost} $\hat{c}_i = c_i + (\Phi(D_i) - \Phi(D_{i-1}))
        = (t+1) + (1 - t) = 2$.
\end{itemize}

The $t$ -- the entire unbounded carry chain -- cancels.
The amortized cost is a flat $2$ no matter how long the carry ripples. The
expensive carries are paid for exactly by the drop in potential they cause: each
cleared bit was a unit of stored potential, released to fund its own flip.

\begin{intuitionbox}
$\Phi$ counts the $1$-bits -- the ``loaded springs'' of the counter. A long carry
chain is a long line of springs releasing at once: lots of actual work, but each
release hands back the potential that was stored when that bit was set. The set
bit at the end is a new spring loaded, costing the one unit that does not cancel.
Amortized: $2$, always.
\end{intuitionbox}

Therefore $n$ increments cost at most $2n + \Phi(D_0) - \Phi(D_n) \le 2n = O(n)$,
matching the aggregate and accounting analyses a third time.

\begin{remark}
What about overflow -- when the counter is all $1$s and \textsc{Increment} wraps it
to all $0$s? Then it clears all $k$ bits and sets none, so $c_i = k$ and
$\Phi$ drops from $k$ to $0$, a change of $-k$. The amortized cost is
$k + (-k) = 0 \le 2$, still bounded. The telescoping bound still holds because
$\Phi(D_n) \ge 0$ regardless. If you forbid overflow (start at $0$, do at most
$2^k - 1$ increments), $\Phi$ never reaches the wrap and the clean ``$\le 2n$''
holds outright.
\end{remark}

\begin{trapbox}
A potential function can be \emph{wrong} -- it just has to satisfy
$\Phi(D_0) = 0$ and $\Phi \ge 0$ to give a \emph{valid} bound, but it gives a
\emph{tight} bound only if you chose it well. Picking $\Phi \equiv 0$ always
satisfies the conditions and makes amortized cost equal actual cost -- useless. The
skill is choosing a $\Phi$ that rises during cheap operations (so their amortized
cost absorbs a little extra) and plunges during expensive ones (so the plunge pays
the bill). If your $\Phi$ does not drop when the expensive operation fires, you
chose the wrong $\Phi$.
\end{trapbox}

\section{The headline example: dynamic tables}

Everything so far has been a warm-up for the example that matters most in
practice. You want an array-backed structure -- a list, a hash table, a stack --
that you can keep inserting into, but you do not know in advance how big it will
get. Arrays have fixed size. So what do you do when the array fills up?

You allocate a bigger one, copy everything over, and keep going. The question is:
\emph{how much bigger?} The answer, and the reason it is the right answer, is one
of the cleanest results amortized analysis delivers.

\subsection{The setup and the naive disaster}

A \emph{dynamic table} stores items in an array of some capacity. It tracks two
numbers: $\mathit{size}$ (how many items are stored) and $\mathit{capacity}$ (how
many the current array can hold). The \emph{load factor} is
$\alpha = \mathit{size} / \mathit{capacity}$. We \textsc{Insert} one item at a
time. When the array is full ($\mathit{size} = \mathit{capacity}$), an insert must
first \emph{expand}: allocate a new array, copy all existing items in, free the
old one, and only then store the new item.

\begin{trapbox}
The tempting, terrible idea: when the table fills, grow it by a \emph{fixed} amount
-- say one extra slot, or a constant $c$ slots. Then \emph{every} insert (or every
$c$th insert) triggers a full copy. If you grow by one slot each time, inserting
$n$ items copies $0 + 1 + 2 + \cdots + (n-1) = \Theta(n^2)$ items total. That is
quadratic total work from what looks like a harmless ``just make it a bit
bigger.'' Growing additively is the classic performance bug.
\end{trapbox}

The fix is to grow \emph{multiplicatively}: when the table fills, \emph{double} its
capacity. This is called \emph{table doubling}. Here is the insert operation.

\begin{algorithm}[H]
\caption{\textsc{Table-Insert}($T, x$)}
\KwIn{A dynamic table $T$ with fields $T.\mathit{size}$, $T.\mathit{cap}$, and array $T.A$.}
\KwOut{$x$ appended; the table is doubled first if it was full.}
\If{$T.\mathit{size} = T.\mathit{cap}$}{
  $\mathit{newcap} \leftarrow \max(1,\; 2 \cdot T.\mathit{cap})$\;
  allocate new array $B$ of capacity $\mathit{newcap}$\;
  \For{$j \leftarrow 1$ \KwTo $T.\mathit{size}$}{
    $B[j] \leftarrow T.A[j]$ \tcp*{copy every existing item: the expensive part}
  }
  free $T.A$;\quad $T.A \leftarrow B$;\quad $T.\mathit{cap} \leftarrow \mathit{newcap}$\;
}
$T.\mathit{size} \leftarrow T.\mathit{size} + 1$\;
$T.A[T.\mathit{size}] \leftarrow x$ \tcp*{store the new item: always cheap}
\end{algorithm}

We charge $1$ unit per item stored or copied. A non-expanding insert costs $1$ (it
just stores $x$). An expanding insert that copies $m$ existing items costs
$m + 1$. Let us watch the cost of the $i$th insert, starting from an empty table.

The expansions happen exactly when the table is full and one more item arrives.
With doubling from capacity $1$, the table expands on inserts numbered
$1, 2, 3, 5, 9, 17, \ldots$. After the first three (which walk the capacity through
$1, 2, 4$), the pattern settles down: for $j \ge 2$, insert number $2^j + 1$ is the
one that fills capacity $2^j$ and forces the doubling to $2^{j+1}$. (Insert $1$
fills capacity $1$; insert $2$ forces capacity to $2$; insert $3$ forces capacity
to $4$; insert $5$ forces capacity to $8$; insert $9$ to $16$; and so on.)

\begin{figure}[h]
\centering
\begin{tikzpicture}[scale=0.62]
  % bars: cost of insert i, for i=1..16
  \def\costs{1/1, 2/2, 3/3, 4/1, 5/5, 6/1, 7/1, 8/1, 9/9, 10/1, 11/1, 12/1, 13/1, 14/1, 15/1, 16/1}
  \foreach \i/\c in {1/1,2/2,3/3,4/1,5/5,6/1,7/1,8/1,9/9,10/1,11/1,12/1,13/1,14/1,15/1,16/1}{
    \fill[sourceblue!70] (\i-0.4,0) rectangle (\i+0.4,\c*0.5);
    \node[font=\tiny] at (\i,-0.45) {\i};
  }
  \draw[->] (0.2,0) -- (17,0) node[right,font=\small] {insert \#};
  \draw[->] (0.2,0) -- (0.2,5.2) node[above,font=\small] {cost};
  \node[font=\tiny,left] at (0.2,0.5) {1};
  \node[font=\tiny,left] at (0.2,2.5) {5};
  \node[font=\tiny,left] at (0.2,4.5) {9};
  \draw[dashed,trapred] (0.2,1) -- (17,1) node[right,font=\tiny,trapred] {amortized $\le 3$};
\end{tikzpicture}
\caption{Cost of each insert into a doubling table, $i = 1 \ldots 16$. Almost
every insert costs $1$. The spikes at $i = 2^j + 1$ (heights $2,3,5,9,\ldots$) are
the doublings. Tall but rare -- and as we will prove, their total is bounded by a
geometric series.}
\label{ch17:doubling-spikes}
\end{figure}

\subsection{Aggregate analysis of doubling}

Consider $n$ inserts into an initially empty table. Each insert pays $1$ for
storing its own item: that is $n$ total. On top of that, the expanding inserts pay
for \emph{copies}. How many copies happen in total?

The $j$th expansion copies $2^{j-1}$ items (it copies everything in the table just
before doubling capacity from $2^{j-1}$ to $2^{j}$). The expansions that occur
during $n$ inserts copy
\[
  1 + 2 + 4 + 8 + \cdots + 2^{\lfloor \log_2(n-1) \rfloor} \;<\; 2n.
\]
Let us be exact rather than hand-wavy. The copies form the geometric sum
$\sum_{j=0}^{J} 2^{j}$ where $2^{J} < n$, and a geometric sum of powers of two is
just under twice its largest term:
\[
  \sum_{j=0}^{J} 2^{j} \;=\; 2^{J+1} - 1 \;<\; 2 \cdot 2^{J} \;<\; 2n.
\]
So total copies $< 2n$, and the total cost is
\[
  T(n) \;=\; \underbrace{n}_{\text{stores}} \;+\; \underbrace{(\text{copies})}_{<\,2n}
  \;<\; n + 2n \;=\; 3n \;=\; O(n).
\]
The amortized cost per insert is at most $3n/n = 3 = O(1)$.

\begin{aestheticbox}
This is the result that makes growable arrays -- C++ \texttt{vector}, Java
\texttt{ArrayList}, Python \texttt{list}, Go slices -- practical. Each one of those
doubles (or grows by a constant factor like $1.5\times$) precisely so that the
amortized cost of appending is $O(1)$. Without amortized analysis you might believe
appending is sometimes $O(n)$ and panic; with it, you can promise users
``\texttt{append} is constant time on average over any sequence,'' and mean it as a
hard guarantee. Once again, the geometric series bounds the total work: all the rare doublings
combined cost at most a constant multiple of the cheap appends.
\end{aestheticbox}

\begin{intuitionbox}
\textbf{The accounting view of table doubling.}
We can also understand the $\$3$ amortized cost directly through credit tokens. When we insert an item into a table of capacity $m$ that was just doubled (so it currently holds $m/2$ elements), we charge an amortized fee of $\$3$:
\begin{itemize}
  \item $\$1$ pays for inserting the element into its empty array slot.
  \item $\$1$ is stored as a credit token on the new element, to pay for moving itself during the next doubling.
  \item $\$1$ is deposited as a credit token onto one of the older elements that was copied in during the previous doubling without credit.
\end{itemize}
By the time the table becomes completely full ($m$ elements), exactly $m/2$ new elements have been inserted, contributing $(m/2) \times \$2 = \$m$ in surplus tokens. Each of the $m$ elements now holds $\$1$ of credit, exactly funding the $m$ copy operations needed to relocate to the new $2m$ array. Figure~\ref{ch17:fig-dynamic-array-accounting} illustrates this accumulation of credit.
\end{intuitionbox}

\begin{figure}[h]
\centering
\begin{tikzpicture}[font=\small]
  % Initial state: capacity 8, 4 items already present (from previous doubling)
  \begin{scope}
    \node[font=\footnotesize\scshape, text=gray!70!black, anchor=west] at (0, 1.8) {Capacity $8$, currently holding $n=4$ items (zero credit right after doubling)};
    \foreach \i in {1,...,8} {
      \ifnum\i<5
        \draw[thick, draw=sourceblue!70, fill=sourceblue!15, rounded corners=1pt] ({(\i-1)*1.3}, 0) rectangle ({(\i-1)*1.3+1.0}, 0.7);
        \node[font=\scriptsize] at ({(\i-1)*1.3+0.5}, 0.35) {Item \i};
      \else
        \draw[thick, draw=gray!40, fill=white, dashed, rounded corners=1pt] ({(\i-1)*1.3}, 0) rectangle ({(\i-1)*1.3+1.0}, 0.7);
        \node[font=\tiny, text=gray!40] at ({(\i-1)*1.3+0.5}, 0.35) {empty};
      \fi
    }
  \end{scope}

  % Transition: inserting items 5, 6, 7, 8 (each paying $3)
  \draw[-{Latex[length=2.5mm]}, very thick, quotientorange] (5.2, -0.4) -- 
    node[right=4pt, font=\scriptsize, text=quotientorange] {Insert items $5, 6, 7, 8$: each pays $\$3$ ($\$1$ insert + $\$2$ credit)} (5.2, -1.4);

  % Full state: capacity 8, 8 items, each carrying $1 token
  \begin{scope}[yshift=-4.0cm]
    \node[font=\footnotesize\scshape, text=gray!70!black, anchor=west] at (0, 1.8) {Table full ($n=8$, $\text{cap}=8$): exactly $\$8$ stored credit, paying for next doubling to $16$};
    \foreach \i in {1,...,8} {
      \draw[thick, draw=targetgreen!80, fill=targetgreen!20, rounded corners=1pt] ({(\i-1)*1.3}, 0) rectangle ({(\i-1)*1.3+1.0}, 0.7);
      \node[font=\scriptsize] at ({(\i-1)*1.3+0.5}, 0.35) {Item \i};
      \node[draw=quotientorange, fill=quotientorange!35, circle, thick, inner sep=1.5pt, minimum size=13pt, font=\tiny\bfseries] 
        at ({(\i-1)*1.3+0.5}, 1.1) {\$1};
    }
    % Annotations showing which insert deposited which coin
    \draw[dotted, thick, quotientorange] (4.4, -0.3) -- (0.5, -0.3) -- (0.5, -0.1);
    \draw[dotted, thick, quotientorange] (4.4, -0.3) -- (4.4, -0.1);
    \node[font=\tiny, text=quotientorange!90!black] at (2.45, -0.55) {Insert \#$5$ banked coins on Item $5$ and Item $1$};
  \end{scope}
\end{tikzpicture}
\caption{The accounting method for dynamic table doubling. Each append pays an amortized charge of $\$3$: $\$1$ performs the immediate insertion into an empty slot, $\$1$ stays on the new item to fund its future relocation, and $\$1$ is deposited onto an earlier item that lacks credit. By the time the table fills from size $m/2$ to $m$, all $m$ items hold $\$1$, generating precisely the $\$m$ credit required to copy all elements into the new $2m$-capacity array.}
\label{ch17:fig-dynamic-array-accounting}
\end{figure}

\subsection{Potential analysis of doubling: choosing $\Phi$}

The aggregate argument is clean, but the potential method is more useful here because it
gives a per-operation account and extends naturally to deletion. We need a $\Phi$
that is $0$ on an empty table, never negative, and that has steadily built up
to a large value exactly when the table is full, so that the expansion can drain
it.

\begin{intuitionbox}
When is a doubling table about to do expensive work? When it is \emph{full}. So we
want potential to be near zero right after a doubling (the table is half empty,
lots of room) and to climb to its peak right before the next doubling (the table is
packed). The potential should measure ``how dangerously full are we?''
\end{intuitionbox}

The classic choice that achieves this:
\[
  \Phi(T) \;=\; 2 \cdot T.\mathit{size} \;-\; T.\mathit{cap}.
\]
Let us sanity-check the two required conditions. Right after a genuine doubling
(capacity at least $2$), the table is half full: $\mathit{size} = \mathit{cap}/2$,
so $\Phi = 2(\mathit{cap}/2) - \mathit{cap} = 0$. As we insert without expanding,
$\mathit{size}$ rises while $\mathit{cap}$ stays fixed, so $\Phi$ climbs. When the
table becomes full, $\mathit{size} = \mathit{cap}$, so
$\Phi = 2\mathit{cap} - \mathit{cap} = \mathit{cap}$ -- a big reservoir, exactly
equal to the number of items the next expansion will have to copy. The potential
has saved up precisely enough.

We also need $\Phi \ge 0$ throughout. Because doubling keeps the load factor
$\alpha = \mathit{size}/\mathit{cap}$ at least $\tfrac12$ at all times after the
first item (right after a double it is exactly $\tfrac12$, and it only rises until
the next double), we have $\mathit{size} \ge \mathit{cap}/2$, hence
$2\,\mathit{size} - \mathit{cap} \ge 0$. And the empty table has $\Phi = 0$. Good --
both conditions hold.

\begin{figure}[h]
\centering
\begin{tikzpicture}[font=\footnotesize,
    full/.style={draw,minimum width=0.45cm,minimum height=0.45cm,fill=sourceblue!30},
    empty/.style={draw,minimum width=0.45cm,minimum height=0.45cm,fill=white}]
  % Show 4 snapshots of the array as it grows. Below each: Phi = 2*size - cap.
  % Snapshot A: cap 4, size 2 (just doubled, half full): Phi = 0
  \begin{scope}[shift={(0,0)}]
    \foreach \i in {0,1}{\node[full] at (0.45*\i,0) {};}
    \foreach \i in {2,3}{\node[empty] at (0.45*\i,0) {};}
    \node[font=\scriptsize,align=center] at (0.675,-0.75) {cap 4, size 2\\$\Phi=0$};
  \end{scope}
  \draw[-{Stealth[length=2.4mm]}] (2.1,0) -- (2.8,0);
  % Snapshot B: cap 4, size 3: Phi = 2
  \begin{scope}[shift={(2.95,0)}]
    \foreach \i in {0,1,2}{\node[full] at (0.45*\i,0) {};}
    \foreach \i in {3}{\node[empty] at (0.45*\i,0) {};}
    \node[font=\scriptsize,align=center] at (0.675,-0.75) {cap 4, size 3\\$\Phi=2$};
  \end{scope}
  \draw[-{Stealth[length=2.4mm]}] (5.05,0) -- (5.75,0);
  % Snapshot C: cap 4, size 4 (full): Phi = 4 (reservoir full)
  \begin{scope}[shift={(5.9,0)}]
    \foreach \i in {0,1,2,3}{\node[full] at (0.45*\i,0) {};}
    \node[font=\scriptsize,align=center] at (0.675,-0.75) {cap 4, size 4\\$\Phi=4$};
  \end{scope}
  \draw[-{Stealth[length=2.4mm]},trapred,thick] (8.0,0) -- (8.7,0)
    node[midway,above,font=\scriptsize,trapred] {double};
  % Snapshot D: cap 8, size 5 (just expanded): Phi = 2
  \begin{scope}[shift={(8.85,0)}]
    \foreach \i in {0,1,2,3,4}{\node[full] at (0.45*\i,0) {};}
    \foreach \i in {5,6,7}{\node[empty] at (0.45*\i,0) {};}
    \node[font=\scriptsize,align=center] at (1.575,-0.75) {cap 8, size 5\\$\Phi=2$};
  \end{scope}
\end{tikzpicture}
\caption{Potential $\Phi = 2\,\mathit{size} - \mathit{cap}$ filling and draining.
Right after a doubling the array is half full and $\Phi$ is small. Each cheap
insert raises $\Phi$ by $2$, building a reservoir. When the array fills,
$\Phi = \mathit{cap}$ -- exactly the number of items the next doubling must copy.
The expensive copy then drains $\Phi$ back down, and the cycle repeats.}
\label{ch17:table-potential}
\end{figure}

Now compute the amortized cost of the $i$th insert. Write $s = T.\mathit{size}$ and
$c = T.\mathit{cap}$ for the values \emph{just before} this insert.
Figure~\ref{ch17:table-potential} previews the arithmetic: watch $\Phi$ climb on
cheap inserts and crash on the doubling.

\textbf{Case 1: no expansion} ($s < c$). Actual cost $c_i = 1$. After insert,
size is $s+1$, capacity unchanged at $c$. So
\[
  \Delta\Phi = \bigl(2(s+1) - c\bigr) - \bigl(2s - c\bigr) = 2.
\]
Amortized cost $\hat{c}_i = 1 + 2 = 3$.

\textbf{Case 2: expansion} ($s = c$, with $c \ge 1$). The insert copies all $c$
items and stores $1$ new one, so actual cost $c_i = c + 1$. Capacity becomes $2c$;
size becomes $c + 1 = s + 1$. The potential goes from
$\Phi_{\text{before}} = 2c - c = c$ to $\Phi_{\text{after}} = 2(c+1) - 2c = 2$. So
\[
  \Delta\Phi = 2 - c.
\]
Amortized cost $\hat{c}_i = (c + 1) + (2 - c) = 3$.

\begin{focusbox}
Both cases give amortized cost exactly $\mathbf{3}$. In Case 2 the $+c$ of
copying is cancelled by the $-c$ drop in potential: the reservoir built up during
the cheap inserts is spent all at once to pay for the copy. Amortized cost per insert is $3 = O(1)$, so $n$
inserts cost $\le 3n = O(n)$, confirming the aggregate result with a tidy
per-operation story.
\end{focusbox}

\begin{remark}
There is one edge case at the very start. The Case~2 algebra assumed $c \ge 1$.
The first insert into a fresh table has $s = c = 0$, and our algorithm grows the
capacity to $\max(1, 2 \cdot 0) = 1$ rather than to $2c = 0$. That insert copies
nothing, so its actual cost is $1$, and the potential moves from $0$ to
$2(1) - 1 = 1$, giving amortized cost $1 + 1 = 2$ -- not $3$. This is harmless: $2$
is still at most $3$, so the bound $\le 3n$ holds for every $n$. From the second
insert onward the clean ``always exactly $3$'' story takes over. (You can see the
first insert costing $2$ in the lab below.)
\end{remark}

\begin{labbox}
Verify the telescoping by hand. Insert $5$ items into an empty doubling table and
track, after each insert: actual cost $c_i$, size, capacity, $\Phi = 2\,\mathit{size}
- \mathit{cap}$, and amortized $\hat{c}_i = c_i + \Delta\Phi$.
You should find: insert 1 (expand to cap 1): $c=1$, $\Phi=1$, $\hat{c}=2$; insert 2
(expand to 2): $c=2$, $\Phi=2$, $\hat c=3$; insert 3 (expand to 4): $c=3$,
$\Phi=2$, $\hat c=3$; insert 4: $c=1$, $\Phi=4$, $\hat c=3$; insert 5 (expand to 8):
$c=5$, $\Phi=2$, $\hat c=3$. Total actual $= 1+2+3+1+5 = 12$. Total amortized
$= 2+3+3+3+3 = 14$. Confirm $12 \le 14$, and that the gap $14 - 12 = 2$ equals the
final $\Phi = 2$ -- exactly the leftover reservoir, as the telescoping formula
predicts.
\end{labbox}

\subsection{Why constant-factor growth, and not additive}

We can now say precisely why doubling -- or any constant-factor growth -- works and
additive growth fails. With multiplicative growth by a factor $r > 1$, the copy
costs across all expansions form a geometric series with ratio $r$, which sums to a
constant times its largest term, giving $O(n)$ total. With additive growth, the
copy costs form an arithmetic series $0 + c + 2c + \cdots$, which sums to
$\Theta(n^2)$.

\begin{complexitybox}
\textbf{Dynamic table, doubling growth (factor $2$):}
\begin{itemize}
  \item Single \textsc{Insert}, worst case: $\Theta(n)$ (when it triggers a copy).
  \item \textsc{Insert}, \emph{amortized}: $\Theta(1)$ (we proved $\le 3$).
  \item $n$ inserts, total: $\Theta(n)$.
  \item Wasted space: at most a factor of $2$ (load factor $\ge \tfrac12$).
\end{itemize}
\textbf{Additive growth (by $c$ each time):} single insert $\Theta(n)$, but
amortized $\Theta(n)$ as well -- $n$ inserts cost $\Theta(n^2)$. Never do this.
\end{complexitybox}

\begin{walkbox}
On your walk, weigh the trade-off. Growing by $2\times$ wastes up to half the
array but appends in amortized $O(1)$. Growing by $1.5\times$ wastes at most a
third and is still amortized $O(1)$ -- the geometric series still converges. Where
is the line? Show yourself that \emph{any} growth factor $r > 1$ gives amortized
$O(1)$, with the constant depending on $r$ as roughly $r/(r-1)$. What happens to
that constant as $r \to 1^{+}$? That limit is exactly why additive growth (morally
$r \to 1$) explodes.
\end{walkbox}

\subsection{A word on deletion and shrinking}

If a table can grow, you may want it to shrink when it empties out, to reclaim
memory. The naive idea -- halve the array as soon as it drops below half full --
contains a trap.

\begin{trapbox}
Suppose you double when full and halve when the table falls to exactly half full.
Now sit right at the boundary $\mathit{size} = \mathit{cap}/2$ and alternate:
insert (table fills, double, copy $\Theta(n)$ items), delete (table drops to half,
halve, copy $\Theta(n)$ items), insert (double again), delete (halve again)\ldots
\emph{Every single operation} now triggers a full copy. The amortized cost rockets
back up to $\Theta(n)$. The structure ``thrashes'' at the threshold.
\end{trapbox}

\begin{figure}[h]
\centering
\begin{tikzpicture}[font=\footnotesize]
  % A horizontal load-factor axis from 0 to 1 for each policy.
  % Policy 1: shrink at 1/2, grow at 1 -> no gap, thrashing.
  \node[font=\small\bfseries] at (-1.6,1.0) {naive};
  \draw[thick] (0,1.0) -- (8,1.0);
  \foreach \x/\lbl in {0/0, 4/$\tfrac12$, 8/1}{
    \draw (\x,0.9) -- (\x,1.1); \node[font=\scriptsize,below=1pt] at (\x,0.9) {\lbl};
  }
  \draw[-{Stealth[length=2mm]},trapred,thick] (4,1.35) -- (4,1.1)
    node[above,font=\scriptsize,trapred] {shrink};
  \draw[-{Stealth[length=2mm]},trapred,thick] (8,1.35) -- (8,1.1)
    node[above,font=\scriptsize,trapred] {grow};
  \draw[decorate,decoration={brace,amplitude=4pt},trapred] (8,0.3) -- (4,0.3)
    node[midway,below=2pt,font=\scriptsize,trapred] {after a resize you land on a trigger};

  % Policy 2: shrink at 1/4, grow at 1, land at 1/2 -> wide gap.
  \node[font=\small\bfseries] at (-1.6,-1.6) {hysteresis};
  \draw[thick] (0,-1.6) -- (8,-1.6);
  \foreach \x/\lbl in {0/0, 2/$\tfrac14$, 4/$\tfrac12$, 8/1}{
    \draw (\x,-1.7) -- (\x,-1.5); \node[font=\scriptsize,below=1pt] at (\x,-1.7) {\lbl};
  }
  \draw[-{Stealth[length=2mm]},targetgreen,thick] (2,-1.25) -- (2,-1.5)
    node[above,font=\scriptsize,targetgreen] {shrink};
  \draw[-{Stealth[length=2mm]},targetgreen,thick] (8,-1.25) -- (8,-1.5)
    node[above,font=\scriptsize,targetgreen] {grow};
  \filldraw[targetgreen] (4,-1.6) circle (2.2pt);
  \draw[decorate,decoration={brace,amplitude=4pt,mirror},targetgreen] (2,-2.3) -- (8,-2.3)
    node[midway,below=2pt,font=\scriptsize,targetgreen!80!black] {land at $1/2$, far from both triggers};
\end{tikzpicture}
\caption{Load-factor thresholds. Top: shrinking at $\tfrac12$ and growing at $1$
leaves the table sitting on the shrink trigger right after a doubling, so a single
insert--delete pair thrashes. Bottom: shrinking at $\tfrac14$ instead leaves every
resize at load factor $\tfrac12$, a comfortable distance from both triggers, so
$\Omega(\mathit{cap})$ operations must pass before the next resize.}
\label{ch17:hysteresis}
\end{figure}

The cure is \emph{hysteresis}: separate the grow and shrink thresholds so the table
must change substantially before either fires (Figure~\ref{ch17:hysteresis}). The standard fix is to double when
full but only halve when the load factor drops to $\tfrac14$ (not $\tfrac12$). After
either resize the table sits at load factor $\tfrac12$, comfortably between the
$\tfrac14$ and $1$ triggers, so $\Omega(\mathit{cap})$ operations must occur before
the next resize. A suitable potential -- something like
$\Phi = |\,2\,\mathit{size} - \mathit{cap}\,|$ when $\alpha \ge \tfrac12$ and
$\Phi = \tfrac12\,\mathit{cap} - \mathit{size}$ when $\alpha < \tfrac12$ -- then
yields amortized $O(1)$ for \emph{both} insert and delete. We leave the careful
case analysis to the exercises and to CLRS, but the idea carries over to other structures.

\begin{focusbox}
Whenever you have a structure that resizes at a threshold, separate the grow and
shrink triggers so a resize leaves you far from \emph{both}. The gap (hysteresis)
guarantees many cheap operations between expensive ones, which is precisely the
condition amortization needs. Resize-at-the-same-threshold is a recurring
real-world bug -- in hash tables, buffers, thread pools.
\end{focusbox}

\section{Choosing the right method}

You now have three methods for the same kind of bound. They are interchangeable in
power -- anything one proves, the others can prove -- but they differ in
convenience.

\begin{itemize}
  \item \textbf{Aggregate} is best when a clean global total $T(n)$ has an obvious
        closed form (a geometric series, a count of bit-flips). It is the least
        flexible: it gives every operation the \emph{same} amortized cost, namely
        $T(n)/n$.
  \item \textbf{Accounting} is best when you can point at a concrete, local unit of
        ``stored work'' -- a coin on each stacked object, a dollar on each set bit
        -- and you want different operations to carry different charges.
  \item \textbf{Potential} is best for the trickiest structures, and it is the
        method that scales: it gives a per-operation bound, it composes when a
        structure supports several operations, and choosing a clever $\Phi$ often
        cracks problems the other two struggle with (Fibonacci heaps, splay trees,
        and union-find by rank all lean on potential functions).
\end{itemize}

\begin{historybox}
Amortized analysis was formalized by Robert E.~Tarjan in his 1985 paper \emph{Amortized Computational Complexity}. Before that paper, theorists and systems programmers used ad-hoc financial bookkeeping or physics analogies without a rigorous unifying framework. Tarjan systematized the three classic approaches -- the aggregate method, the accounting method, and the potential method -- drawing directly upon physical analogies of stored potential energy from classical mechanics. Together with Daniel Sleator, Tarjan used these techniques to analyze splay trees, pairing heaps, and the inverse-Ackermann bounds of disjoint sets.

In the 1970s, strict worst-case analysis was the dominant paradigm in computer science. Under that lens, a resizable array (such as C++ \lstinline|std::vector| or Python list) was judged as a slow $\Theta(n)$ operation because doubling the table requires copying all elements. This pessimistic assessment contradicted actual practice, where dynamic arrays routinely outperformed linked lists. The key insight of amortized analysis was that worst-case per-operation bounds are often misleading: an expensive operation can only occur if preceded by a long sequence of cheap operations that deposited sufficient credit or potential to fund it.

Amortized analysis provided a rigorous mathematical foundation for modern memory allocators, dynamic tables, splay trees, and Fibonacci heaps. It allows systems designers to build structures that occasionally perform massive housekeeping without violating sequence-wide $O(1)$ efficiency. The trade-off is that amortized guarantees are sequence-wide averages and give no latency bound for individual operations. In hard real-time systems (such as medical equipment, automotive engine controls, or real-time audio synthesis), a single 10ms table reallocation pause can cause catastrophic buffer underflows or safety violations. Designing de-amortized algorithms that spread expensive resizing work incrementally across individual steps without sacrificing asymptotic bounds remains an active area of system and algorithm design.
\end{historybox}

\begin{aestheticbox}
Why does a working programmer care? Because amortized bounds are the only truthful
way to advertise the performance of growable, self-adjusting structures, and those
are everywhere: every dynamic array, every hash table that rehashes, every rope or
gap buffer in a text editor, every memory allocator that occasionally compacts.
``Worst case per call'' would force you to quote terrifying, useless numbers.
Amortized analysis lets you make a true, strong promise -- ``constant time over any
sequence'' -- and back it with a proof that no adversary can break.
\end{aestheticbox}

With amortized analysis in hand, we complete Part~IV on algorithmic design and analysis
techniques. Throughout our study of sorting, dynamic programming, and greedy choices, we
have quietly assumed a uniform-cost random access machine (RAM) model: reading any word in
memory costs the same unit of time.

In Part~V (Advanced Data Structures), that assumption meets the reality of physical hardware.
When datasets grow too large to fit in main memory, storage shifts to magnetic disks or solid-state
drives, where fetching a block of data takes thousands or millions of CPU clock cycles. In
\textbf{Chapter~18 (B-Trees)}, we meet multi-way balanced search trees whose large branching
factors match the page sizes of secondary storage, minimizing disk accesses while maintaining
logarithmic dictionary operations.

\section{Exercises}

\begin{exercise}[Multipop by hand]
Starting from an empty stack, execute: \textsc{Push}(1), \textsc{Push}(2),
\textsc{Push}(3), \textsc{Push}(4), \textsc{Multipop}($3$), \textsc{Push}(5),
\textsc{Multipop}($10$). Write the stack contents and actual cost after each
operation, sum the actual costs, and verify the total is at most $2$ times the
number of operations. Why did the final \textsc{Multipop}($10$) cost only $2$?
(Warmup.)
\end{exercise}

\begin{exercise}[Aggregate the counter exactly]
For a $5$-bit counter starting at $0$, list the cost (number of bit-flips) of each
of the first $16$ increments. Sum them, and check your sum against the formula
$\sum_{i\ge 0}\lfloor 16/2^i\rfloor$. By what factor does the true total beat the
naive bound $16 \cdot 5$? (Warmup.)
\end{exercise}

\begin{exercise}[Accounting charges must stay solvent]
A student proposes charging \textsc{Push} a flat $\$1$ and \textsc{Pop} a flat
$\$1$, with \textsc{Multipop} charged $\$1$ total (not per item). Exhibit a short
sequence of operations on which this charging scheme lets the ``bank balance'' go
negative -- i.e.\ total charged falls below total real cost -- thereby proving the
scheme invalid. (Medium.)
\end{exercise}

\begin{exercise}[A reset operation]
Augment the binary counter with a \textsc{Reset} operation that sets the entire
counter back to $0$ (cost equal to the number of $1$-bits cleared). Using the
potential $\Phi = (\text{number of }1\text{-bits})$, show that a sequence of $n$
operations, each an \textsc{Increment} or a \textsc{Reset}, still costs $O(n)$.
Identify the amortized cost of \textsc{Reset} under this potential. (Medium.)
\end{exercise}

\begin{exercise}[Telescoping leftover]
Prove, directly from the telescoping identity
$\sum c_i = \sum \hat{c}_i - \Phi(D_n) + \Phi(D_0)$, that for the doubling table
the difference between total amortized and total actual cost after $n$ inserts
equals exactly $\Phi(D_n) = 2\,\mathit{size} - \mathit{cap}$ of the final table.
Verify this against the $5$-insert lab earlier in the chapter, where the gap was
$2$. (Medium.)
\end{exercise}

\begin{exercise}[Tripling instead of doubling]
Suppose the dynamic table \emph{triples} its capacity when full instead of
doubling. Redo the aggregate analysis: show the total copy cost over $n$ inserts is
still $O(n)$, and find the constant (the amortized cost per insert). Then find a
potential function $\Phi$ analogous to $2\,\mathit{size} - \mathit{cap}$ that
proves the same bound per operation. (Medium-to-hard.)
\end{exercise}

\begin{exercise}[General growth factor]
Show that growing the table by a factor $r > 1$ on each expansion gives amortized
insert cost at most $1 + \frac{r}{r-1}$. Evaluate this for $r = 2$, $r = 1.5$, and
$r = 1.1$. Explain, using your formula, why the amortized cost diverges as
$r \to 1^{+}$, connecting this to why additive growth is $\Theta(n^2)$. (Hard.)
\end{exercise}

\begin{exercise}[Why halving at $\tfrac12$ thrashes]
Consider a table that doubles when full and halves the moment size drops to
$\mathit{cap}/2$. Construct an explicit sequence of $n$ insert/delete operations on
which \emph{every} operation triggers a full copy, so the total cost is
$\Theta(n^2)$ and the amortized cost is $\Theta(n)$. Then explain precisely how
moving the shrink threshold to $\mathit{cap}/4$ breaks your construction. (Hard.)
\end{exercise}

\begin{exercise}[Insert and delete with $\tfrac14$ hysteresis]
For the dynamic table that doubles when full and halves when the load factor drops
to $\tfrac14$, consider the potential
\[
  \Phi(T) = \begin{cases}
    2\,\mathit{size} - \mathit{cap}, & \alpha \ge \tfrac12,\\[2pt]
    \tfrac12\,\mathit{cap} - \mathit{size}, & \alpha < \tfrac12.
  \end{cases}
\]
Verify $\Phi \ge 0$ and $\Phi = 0$ on an empty table, then show that both
\textsc{Insert} and \textsc{Delete} have amortized cost $O(1)$ by checking the four
cases: insert with/without expansion, delete with/without contraction. (Challenge.)
\end{exercise}

\begin{exercise}[Designing a potential from scratch]
You maintain a sorted dynamic structure that occasionally performs a global
``rebuild'' costing $\Theta(m)$ when $m$ items have accumulated since the last
rebuild, and rebuilds happen every time the count of pending items reaches the
current $\mathit{size}$. Propose a potential function and prove the amortized cost
per insert is $O(1)$. (This is the seed of the ``rebuild when you double'' idea
behind many real structures.) (Challenge.)
\end{exercise}
